\documentclass[10pt,a4paper]{amsart}
\usepackage[margin=19mm]{geometry}
\usepackage{amsmath,amssymb,mathtools}
\usepackage[scr=rsfs,cal=euler]{mathalfa}
\usepackage{xfrac}
\usepackage[unicode,hidelinks,bookmarks=false]{hyperref}
\newcommand{\ie}{i.e.\ }
\newcommand{\cf}{cf.\ }
\newcommand{\resp}{respectively\ }
\newcommand{\Subsection}{\subsection}
\providecommand{\nobreakdash}{\nobreak\hbox{-}}
\setlength{\emergencystretch}{2em}
\multlinegap0pt
\usepackage{amssymb}
\usepackage{mathtools}
\usepackage[shortlabels]{enumitem}
\newtheorem{theorem}{Theorem}
\newtheorem{lemma}{Lemma}
\newtheorem{claim}{Claim}
\title{Above and below}
\author{Wenchong Chen, Cosmin Pohoata}
\date{Source: arXiv:2605.27061v1 (CC BY 4.0)}
\begin{document}
\maketitle

\noindent\emph{Source and licence.} Above and below. Version arXiv:2605.27061v1 (CC BY 4.0), distributed by arXiv under the Creative Commons Attribution 4.0 International licence, http://creativecommons.org/licenses/by/4.0/, as recorded in the arXiv licence metadata for this exact version and shown on the abstract page https://arxiv.org/abs/2605.27061v1. Full licence notice: \texttt{licenses/arxiv-2605.27061-CC-BY-4.0.txt}.\par\smallskip

\noindent\textit{Editorial scope.} Counted unit: the article's theorem thm:higher-dimensional (source0.tex lines 148--159). The article's complete original proof of it is delivered with it (source0.tex lines 451--591, one \texttt{proof} environment of 141 lines, which the article places away from the statement). No mathematics is written, repaired or re-derived: the statement and the proof are the article's own lines, in the article's own notation, and no retained line is substituted or re-flowed.  Retained uncounted as the source-local context the proof needs: lines 269--280; lines 333--339; lines 359--372.  Nothing was omitted from any retained span, so the package carries no omission record.  No retained line contains a citation, so the wrapper carries no bibliography and no bibliography entry.  No retained line contains a figure environment, an image-inclusion command or drawing code, so no image is delivered and the package declares no asset dependency.  Editorial formatting only, in addition: an AMS \texttt{amsart} preamble that restates the article's own declarations and macros used by the retained lines, namely the environments \texttt{theorem}, \texttt{lemma}, \texttt{claim} on separate counters, together with the standard packages those lines need.  The retained environments print in this document as Theorem 1, Lemma 1, Claim 1 in order of appearance, whereas the article prints the same items with the numbering of its own full document.  The complete native source of the article as distributed in the arXiv e-print for this exact version is bundled unchanged as \texttt{sources/201466/source0.tex}.
\begin{lemma}\label{lem:color-det}
For a sequence $x_0,\ldots,x_d\in\mathbb R^d$ with cyclic projections, its
above--below color is positive if and only if
\[
\det
\begin{pmatrix}
1&1&\cdots&1\\
x_0&x_1&\cdots&x_d
\end{pmatrix}
>0.
\]
\end{lemma}

\begin{lemma}\label{lem:cramer}
Let $P$ be an $n\times(n+1)$ matrix. For $0\le j\le n$, let $P_j$ denote the
$n\times n$ matrix obtained by deleting the $(j+1)$-st column of $P$. Then
\[
\bigl(\det P_0,-\det P_1,\ldots,(-1)^n\det P_n\bigr)^T\in \ker P.
\]
\end{lemma}

\begin{lemma}\label{lem:Plucker}
Let \(Q\) be an \(n\times(n+2)\) matrix, with columns $v_0,v_1,\ldots,v_{n+1}$, where \(n\geq 2\). For \(0\leq i<j\leq n+1\), let $\delta_{i,j}
  =
  \det(v_0,\ldots,\widehat v_i,\ldots,\widehat v_j,\ldots,v_{n+1})$, where the remaining \(n\) columns are kept in their original order. Then, for every $0\leq i_1<i_2<i_3<i_4\leq n+1$, we have
\begin{equation}\label{eq:plucker}
  \delta_{i_1,i_2}\delta_{i_3,i_4}
  -
  \delta_{i_1,i_3}\delta_{i_2,i_4}
  +
  \delta_{i_1,i_4}\delta_{i_2,i_3}
  =
  0.
\end{equation}
\end{lemma}

\begin{theorem}\label{thm:higher-dimensional}
For every fixed $d\ge 2$,
\[
\operatorname{EM}^{(d)}(k)
=
\operatorname{AB}^{(d)}_{\mathrm{mc}}(k)
\le
\operatorname{AB}^{(d)}(k)
\le
\overline R_{\mathrm{mon}}(k;d+1).
\]
\end{theorem}

\begin{proof}[Proof of Theorem~\ref{thm:higher-dimensional}]
We first prove that the above--below coloring is $(d+1)$-monotone. It is enough to prove the one-switch property on every ordered $(d+2)$-tuple. Let $x_0,x_1,\ldots,x_{d+1}\in\mathbb R^d$ have cyclic projections, and write $x_i=(z_i,h_i)$, where
$z_i\in\mathbb R^{d-1}$. Put
\[
P=
\begin{pmatrix}
1&1&\cdots&1\\
x_0&x_1&\cdots&x_{d+1}
\end{pmatrix},
\qquad
Q=
\begin{pmatrix}
1&1&\cdots&1\\
z_0&z_1&\cdots&z_{d+1}
\end{pmatrix}.
\]
Thus $P$ is a $(d+1)\times(d+2)$ matrix and $Q$ is the $d\times(d+2)$ matrix
obtained from $P$ by deleting the last row.

For $0\le j\le d+1$, let $D_j$ be the determinant of the
$(d+1)\times(d+1)$ matrix obtained from $P$ by deleting the $(j+1)$-st column.
By Lemma~\ref{lem:color-det}, the above--below color of the $(d+1)$-tuple $\{x_0,\ldots,x_{d+1}\}\setminus\{x_j\}$ is the sign of $D_j$.

For $0\le a<b\le d+1$, let $\delta_{a,b}$ be the determinant of the
$d\times d$ matrix obtained from $Q$ by deleting the $(a+1)$-st and $(b+1)$-st
columns. By the cyclic projection condition, and with the same orientation
convention as above, all these minors are positive: $\delta_{a,b}>0$.

By Lemma~\ref{lem:cramer}, applied to $P$, the vector $\vec v_0=(D_0,-D_1,\ldots,(-1)^dD_d,(-1)^{d+1}D_{d+1})^T$ lies in $\ker P$, and hence also in $\ker Q$. Applying Lemma~\ref{lem:cramer}
to the first $d+1$ columns of $Q$ and to the last $d+1$ columns of $Q$, and
then extending by a zero coordinate, gives
\[
\vec v_1=(\delta_{0,d+1},-\delta_{1,d+1},\ldots,(-1)^d\delta_{d,d+1},0)^T,\ \ \vec v_2=(0,-\delta_{0,1},\delta_{0,2},\ldots,(-1)^{d+1}\delta_{0,d+1})^T
\in\ker Q.
\]
Since $Q$ has rank $d$, its kernel is two-dimensional, and $\vec v_1,\vec v_2$
form a basis. Hence $\vec v_0=\alpha\vec v_1+\beta\vec v_2$ for some real
numbers $\alpha,\beta$. Comparing the first and last coordinates gives
\[
\alpha=\frac{D_0}{\delta_{0,d+1}},
\qquad
\beta=\frac{D_{d+1}}{\delta_{0,d+1}}.
\]
Therefore, for every $1\le j\le d$,
\begin{equation}\label{eq:Dj-monotone-combination}
D_j
=
\frac{D_0\,\delta_{j,d+1}+D_{d+1}\,\delta_{0,j}}
     {\delta_{0,d+1}}.
\end{equation}

We now use the Pl\"ucker relation from Lemma \ref{lem:Plucker} for $(i_1,i_2,i_3,i_4)=(0,j,j+1,d+1)$. We get
\begin{equation}\label{eq:plucker-monotone-ratio}
\delta_{0,j+1}\delta_{j,d+1}
=
\delta_{0,j}\delta_{j+1,d+1}
+
\delta_{0,d+1}\delta_{j,j+1},
\qquad 1\le j\le d-1.
\end{equation}
Since every term in \eqref{eq:plucker-monotone-ratio} is positive, note that this implies that the ratios
\[
\rho_j:=\frac{\delta_{j,d+1}}{\delta_{0,j}},
\qquad 1\le j\le d,
\]
are strictly decreasing. This can be leveraged as follows. If $D_0$ and $D_{d+1}$ have the same sign, then
\eqref{eq:Dj-monotone-combination} implies that every $D_j$, $1\le j\le d$,
has this same sign. Suppose instead that $D_0$ and $D_{d+1}$ have opposite
signs. If $D_0>0$ and $D_{d+1}<0$, then
\eqref{eq:Dj-monotone-combination} gives
\[
D_j>0
\quad\Longleftrightarrow\quad
\rho_j>-\frac{D_{d+1}}{D_0}.
\]
However, we know that $\rho_1>\rho_2>\cdots>\rho_d$, thereby the signs of $D_0,D_1,\ldots,D_d,D_{d+1}$ change at most once. The case $D_0<0$ and $D_{d+1}>0$ is identical, with the
inequality reversed.

Thus the signs of the $D_j$'s form a one-switch sequence in deletion order.
The lexicographic order on the $(d+1)$-subsets of
$\{0,1,\ldots,d+1\}$ is the reverse of the deletion order
$j=0,1,\ldots,d+1$, and reversing a sequence preserves the property of having
at most one sign change. Hence the above--below coloring is
$(d+1)$-monotone. Consequently, $\operatorname{AB}^{(d)}(k)\le \overline R_{\mathrm{mon}}(k;d+1)$.

It remains to prove the moment-curve identity $\operatorname{AB}^{(d)}_{\mathrm{mc}}(k)=\operatorname{EM}^{(d)}(k)$. The main idea is the following generalized Vandermonde identity.

\begin{claim}\label{claim:vandermonde-divdiff}
Let $p_0=(t_0,h_0),\dots,p_d=(t_d,h_d)$, where $t_0<\cdots<t_d$. Then
\[
\det
\begin{pmatrix}
1&1&\cdots&1\\
t_0&t_1&\cdots&t_d\\
t_0^2&t_1^2&\cdots&t_d^2\\
\vdots&\vdots&&\vdots\\
t_0^{d-1}&t_1^{d-1}&\cdots&t_d^{d-1}\\
h_0&h_1&\cdots&h_d
\end{pmatrix}
=
\left(\prod_{0\le i<j\le d}(t_j-t_i)\right)
\Delta_d(p_0,\dots,p_d).
\]
\end{claim}

While this is likely fairly standard as well, we would like to include an argument which we think best explains the connection between higher order divided differences and cyclic projections (this is also the starting point of this paper, in some sense). Let $f(x)=a_0+a_1x+\cdots+a_dx^d$ be the unique polynomial of degree at most $d$ satisfying $f(t_i)=h_i$ for all
$i=0,\dots,d$. Using the Lagrange interpolation formula, it is easy to see that
the leading coefficient $a_d$ is precisely the divided difference $a_d=\Delta_d(p_0,\dots,p_d)$. On the other hand, in the determinant on the left hand side of
Claim~\ref{claim:vandermonde-divdiff}, the last row is
\[
(h_0,\dots,h_d)
=
a_0(1,\dots,1)+a_1(t_0,\dots,t_d)+\cdots+a_d(t_0^d,\dots,t_d^d).
\]
Subtracting from the last row the corresponding linear combination of the
previous rows leaves $a_d(t_0^d,\dots,t_d^d)$ as the last row. Therefore the
determinant becomes
\[
a_d
\det
\begin{pmatrix}
1&1&\cdots&1\\
t_0&t_1&\cdots&t_d\\
t_0^2&t_1^2&\cdots&t_d^2\\
\vdots&\vdots&&\vdots\\
t_0^{d-1}&t_1^{d-1}&\cdots&t_d^{d-1}\\
t_0^d&t_1^d&\cdots&t_d^d
\end{pmatrix}.
\]
The remaining determinant is the usual Vandermonde determinant,
$\prod_{0\le i<j\le d}(t_j-t_i)$, and since
$a_d=\Delta_d(p_0,\dots,p_d)$, this proves Claim~\ref{claim:vandermonde-divdiff}.

Since the Vandermonde factor $\prod_{0\le i<j\le d}(t_j-t_i)$ is positive,
Lemma~\ref{lem:color-det} shows that the above--below color of
$x_0,\ldots,x_d$ in the moment-curve model is determined, up to the fixed sign
convention in the definition of positivity, by the sign of
$\Delta_d(p_0,\ldots,p_d)$, and vice versa. Hence monochromatic subsequences in
the moment-curve model are exactly $d$-th order monotone subsequences. Thus $\operatorname{AB}^{(d)}_{\mathrm{mc}}(k)=\operatorname{EM}^{(d)}(k)$. Since the moment-curve model is a special case of cyclic projections, we also
have $\operatorname{AB}^{(d)}_{\mathrm{mc}}(k)\le \operatorname{AB}^{(d)}(k)$. Together with the monotone-coloring bound proved above, this completes the proof of Theorem~\ref{thm:higher-dimensional}.
\end{proof}

\end{document}
